% ---
% id: reductions/prop:obstruction
% kind: prop
% title: "Obstruction to encoding event bounds in classical targets"
% summary: >-
%   Some three-outcome qubit POVM with a lower bound on the sum of two of its
%   elements has no squashing to a classical target that keeps that bound as an
%   operator inequality. A device fact about a union of outcomes therefore
%   travels as a constraint entry rather than inside a flag target.
% prominence: standard
% status: accepted
% uses: [def:squashing]
% assumes: []
% refs: []
% provenance: newly-derived
% lean:
% ---

The \tref[weight-preserving target]{weight-preserving flag-state target} of the flag-state submodule encodes
one certified weight \emph{per outcome}. For an \emph{event}
$\mathcal{O}_{*}\subseteq\Vset$, a union of outcomes that must not be merged
(for instance because they carry key information), a device may certify only
the sum, $\sum_{v\in\mathcal{O}_{*}}M_v\succeq p_{*}\idon{\Bsys}$. The
following proposition shows that no classical target can absorb such a
bound, whatever the construction.

\begin{proposition}
\label{prop:obstruction}
Call a \tref{target POVM} $F$ on $\Qsys$ \emph{classical} if every $F_v$ is diagonal
in one fixed orthonormal basis $\{\ket{f}\}_f$ of $\Hs{\Qsys}$. There exist a
POVM $M$ on $\Hs{\Bsys}=\mathbb{C}^{2}$, an event
$\mathcal{O}_{*}\subseteq\Vset$ and a constant $p_{*}>0$ with
$\sum_{v\in\mathcal{O}_{*}}M_v\succeq p_{*}\idon{\Bsys}$ such that \emph{no}
pair $(F,\Lambda)$ exists with $F$ classical, $M\dsq_{\Lambda}F$ and
$\sum_{v\in\mathcal{O}_{*}}F_v\succeq p_{*}\idon{\Qsys}$. Concretely,
\begin{equation}\label{eq:obstruction}
  M_1=\tfrac12\dyad{0},\quad
  M_2=\tfrac12\dyad{+},\quad
  M_3=\idon{\Bsys}-M_1-M_2,\quad
  \mathcal{O}_{*}=\{1,2\},\quad
  p_{*}=\tfrac12\Big(1-\tfrac{1}{\sqrt2}\Big),
\end{equation}
with $\ket{+}=(\ket{0}+\ket{1})/\sqrt2$.
\end{proposition}

\begin{proof}
\emph{The data are admissible.} For two rank-one projections with overlap
$|\innerp{0}{+}|=1/\sqrt2$, the operator $\dyad{0}+\dyad{+}$ has trace $2$ and
eigenvalues $1\pm1/\sqrt2$; hence
$\lmin(M_1+M_2)=\tfrac12(1-1/\sqrt2)=p_{*}>0$ and
$\lmax(M_1+M_2)=\tfrac12(1+1/\sqrt2)<1$, so $M_3\succeq0$ and $\{M_1,M_2,M_3\}$
is a POVM with $M_1+M_2\succeq p_{*}\idon{\Bsys}$.

\emph{No classical target encodes the bound.} Suppose $(F,\Lambda)$ existed.
Put $T_f\coloneqq\Lambda^{\dagger}(\dyad{f})$. Then $\{T_f\}_f$ is a POVM on
$\Bsys$: each $T_f\succeq0$ because $\Lambda^{\dagger}$ is positive, and
$\sum_fT_f=\Lambda^{\dagger}(\idon{\Qsys})=\idon{\Bsys}$ by unitality. Writing
$F_v=\sum_fF_v(f)\dyad{f}$ with $F_v(f)\ge0$ and $\sum_vF_v(f)=1$ (POVM
completeness read on the diagonal), linearity gives
\begin{equation*}
  M_v=\Lambda^{\dagger}(F_v)=\sum_fF_v(f)\,T_f .
\end{equation*}
The encoded event bound $\sum_{v\in\mathcal{O}_{*}}F_v\succeq p_{*}\idon{\Qsys}$
reads, on the diagonal, $F_1(f)+F_2(f)\ge p_{*}>0$ for every $f$, so every $f$
has $F_1(f)>0$ or $F_2(f)>0$.

Suppose $F_1(f)>0$. Every summand of $M_1=\sum_{f'}F_1(f')T_{f'}$ is positive
semidefinite, so $0\preceq F_1(f)T_f\preceq M_1=\tfrac12\dyad{0}$. A positive
operator dominated by a rank-one positive operator is a non-negative multiple
of it: if $X\preceq c\dyad{0}$ with $X\succeq0$, then for every
$\ket{\psi}\perp\ket{0}$ we get $\bra{\psi}X\ket{\psi}\le0$, hence
$\bra{\psi}X\ket{\psi}=0$ and $X\ket{\psi}=0$, so $\operatorname{supp}X
\subseteq\operatorname{span}\ket{0}$. Thus $T_f\propto\dyad{0}$. Symmetrically
$F_2(f)>0$ forces $T_f\propto\dyad{+}$. A non-zero operator cannot be
proportional to both, since $\dyad{0}$ and $\dyad{+}$ are distinct rank-one
projections; so $F_1(f)F_2(f)>0$ forces $T_f=0$, and in every case each $T_f$
is $0$, or a non-negative multiple of $\dyad{0}$, or a non-negative multiple of
$\dyad{+}$.

Summing, $\idon{\Bsys}=\sum_fT_f=\alpha\dyad{0}+\beta\dyad{+}$ for some
$\alpha,\beta\ge0$ (two non-negative scalars local to this proof, unrelated
to the R\'enyi order $\alpha$). Testing against $\ket{1}$ gives
$1=\bra{1}\big(\alpha\dyad{0}+\beta\dyad{+}\big)\ket{1}=\beta/2$, so
$\beta=2$; testing the off-diagonal entry gives
$0=\bra{0}\big(\alpha\dyad{0}+\beta\dyad{+}\big)\ket{1}=\beta/2$, so
$\beta=0$ --- a contradiction.
\end{proof}
